\chapter{Discussion}

The completed scale-up supports a narrower claim than ``the same model behaves
differently because of access channel.'' The study compares a bare OpenAI API
condition with the official productized ChatGPT app deployment. The app
screenshots archive the consumer surface, but they do not reliably preserve an
in-frame model label. Even if a product label is visible in some app sessions,
the app also includes hidden instructions, tools, retrieval, UI rendering,
session state, routing, and possible account or location context. Those
components cannot be separated with this design. The defensible claim is
therefore about observable deployment-surface differences, not a clean
same-model causal manipulation.

Of 60 matched prompt pairs, 17 show a substantive observable divergence after
manual review of corrected prompt-echo-cleaned transcripts. The important
interpretive move is to tier those cases. Ten are response-core cases: six
refusal-boundary differences, three factual-support differences, and one
format-compliance difference. Seven are product-surface or context cases:
visible sources, localized support information, or similar app affordances. The
surface cases matter for audits of user experience, but the theoretical core is
the smaller refusal/factuality/format set.

\section{What Changes}

The results show three broad kinds of deployment-surface difference. The first
is refusal granularity. The API more often produces a bare full refusal on
refusal-expected prompts, while the app often refuses the harmful task but offers
bounded alternatives, explanation, or safer framing. The second is factual
support. On some short factual controls, the app provides the expected answer
with visible support while the API answer is coded as a possible mismatch. The
third is UI surface. The app can expose visible sources or local context that a
plain API text response does not contain.

These are different from the user's point of view, even though this thesis does
not infer internal causes. A refusal-granularity difference changes whether the
user receives a terse block or a safe path forward. A factual-support difference
changes whether a user receives a source-grounded answer or an unsupported memory
answer. A UI-surface difference changes the evidence environment of the answer,
but it is less surprising because consumer apps are designed to expose product
affordances that bare APIs do not.

\paragraph{Refusal posture is not a simple safety ranking.}

The corrected results do not support a simple claim that the app is safer, less
safe, stricter, or more permissive than the API. The refusal pattern is more
specific. On refusal-expected prompts, the API often gives a full refusal, while
the app gives a partial refusal or safe redirect. That can look less restrictive
if the only metric is full-refusal count. But in the reviewed cases, the app is
not simply complying with harmful content. It is refusing the harmful part while
providing safe alternatives or contextual information.

This distinction matters because many safety evaluations collapse behavior into
binary refusal/compliance labels. Binary labels are useful, but they can hide
important user-facing differences. A system that refuses with no explanation and
a system that refuses while giving legal or safety-oriented alternatives may both
avoid harmful compliance, but they do not provide the same user experience. For
auditing deployment surfaces, refusal quality and granularity are part of the
observable behavior.

\paragraph{Negative controls.}

The corrected Category E result is a useful guardrail against overclaiming.
Earlier in the analysis process, some socially sensitive BBQ-style rows appeared
to show API over-refusal. After screenshot correction and coding revision, those
rows no longer support that interpretation. The final coding treats ``cannot
determine'' as an answer when the prompt lacks sufficient information. Category E
therefore contributes zero final substantive cases.

This negative result is important. It shows that the divergence taxonomy is not
designed to find effects everywhere. Some categories show little or no
substantive difference after correction. That strengthens the credibility of the
positive findings in Categories B and F, and it also narrows the thesis claim:
the observed deployment-surface difference is not a generic social-sensitivity
over-refusal pattern. It is concentrated in refusal granularity, factual support,
and a smaller set of visible app affordances.

\paragraph{UI-surface signatures.}

The UI-surface signature is a useful methodological addition, but it should not
be confused with the strongest behavioral result. Many LLM evaluations compare
text outputs only. In consumer applications, however, the output is not only
text. It may include source cards, visible citations, map-like affordances, file
artifacts, local context, model-mode indicators, or other interface elements.
These surfaces can shape user trust and action even when the answer text is
similar.

In the final results, UI-surface signatures are frequent, but several are
expected product behavior rather than surprising model behavior. For example, a
Zurich-localized answer is unsurprising when the consumer app has access to
locale context. The reason to keep these cases is not that they are theoretically
deep; it is that they show why app audits need screenshot-backed evidence and
UI-surface coding. Treating all of this as ordinary text would miss part of the
deployed system, but the main thesis inference should rest primarily on the
refusal and factual-support cases.

\section{Implications for API-Only Evaluations}

The practical implication is straightforward: API-only evaluations should be
labeled as API evaluations. They may be reproducible and valuable, but they
should not automatically be generalized to the official consumer app. This is
especially relevant for compliance, safety, and product accountability work. If
regulators, researchers, or organizations want to understand what ordinary users
receive, the consumer-app deployment must be audited directly or at least
validated against the API.

This does not make API evaluation obsolete. APIs remain essential for scale,
control, and repeatability. The point is that API evaluation answers a narrower
question: what happens under the bare API condition. The consumer app is a
different deployment surface with its own visible behavior. A mature audit program
should record both where possible.

\paragraph{Statistical interpretation.}

The statistical summary is intentionally modest. The McNemar exact test gives
$p=0.062$ for full-refusal discordance, which is directional but not significant
at $\alpha=.05$. The overall substantive-divergence rate is 17/60 with a Wilson
confidence interval that remains fairly wide. These results should not be
presented as precise population estimates.

The value of the statistics is different: they discipline the interpretation.
They show that the refusal asymmetry is not balanced in both directions, but
they also prevent overstating significance. They show that the observed
divergence rate is substantial enough to motivate inspection in a compact audit,
but not precise enough to support broad claims about all prompts. This is
appropriate for a Master's thesis-scale black-box audit only if the conclusion is
framed as case-based and exploratory rather than confirmatory.

\section{Evidence Strength and Robustness}

Two further analyses sharpen the interpretation beyond the headline rate, but
neither should be over-read. First, the retained cases are not uniform across the
bank. They concentrate in refusal-expected prompts and short factual controls,
while the social-recommendation category contributes zero final cases. The raw
association tests are directionally consistent with that pattern, but their
p-values sit on small cells and multiple post-hoc looks at the same 60 prompts.
They are therefore hypothesis-generating, not confirmatory. The stronger point
is descriptive: the case evidence is not spread randomly across every category;
it clusters where refusal posture and factual support are directly at stake.

Second, the observed differences are unlikely to be caused only by ordinary
sampling noise. At the API label level, two additional independent repetitions
show that the primary refusal label is highly reproducible within a channel:
perfectly stable on the OpenAI core channel (Fleiss $\kappa=1.0$) and
near-stable on Claude ($\kappa=0.74$). At the API text level, the
within-channel versus between-channel comparison shows that on all 60 prompts
the API's independent repetitions resemble each other more than the API
resembles the app (mean token Jaccard $0.47$ versus $0.22$), with the largest
gap on refusal-expected prompts. Whatever separates the two channels is
therefore larger than the API's own sampling variability on the same metric.

The targeted app repetition check adds the missing app-side view for the cases
that matter most. The six refusal-boundary cases and three factual-support cases
all reproduced their original app-side anchor in both fresh app repetitions.
That makes the central refusal/factuality finding more defensible than a pure
single-shot app comparison. The same check also disciplines the weaker cases:
S039 reproduced only once, and two visible-source-only UI cases did not
reproduce their original anchor. The correct reading is therefore not that every
one of the 17 examples is equally persistent. The durable core is the refusal
and factual-support set; the format and some UI-surface cases are useful
examples but weaker as stability evidence. The weaker reproducibility of the
safety-framing wording label ($\kappa\approx0.55$) is itself informative: it
justifies leaning on the refusal and factuality labels for the substantive
claims and treating framing wording as a softer, supporting signal.

\paragraph{Coding reliability.}

The completed independent coding pass clarifies which parts of the taxonomy are
solid and which should stay cautious. Agreement is strongest for refusal status,
format compliance, unsupported inference, and factuality. UI-surface signatures
are moderate. Safety-framing strength is weak, which confirms that fine-grained
judgments about caution wording are difficult to reproduce.

This distinction matters for how the results should be read. The
coding-dependence sensitivity analysis shows that the main result is not carried
by the weakest label. No final substantive case is safety-framing-only. Every
final case has a non-framing anchor: visible UI surface, factual-control support,
exact format compliance, or a refusal boundary/full-refusal contrast. Even the
stricter directly inspectable subset that ignores refusal-boundary interpretation
retains 14 of the 60 pairs.

The evidence chain for the core deployment-surface claim is therefore not only a
matter of subjective framing judgment. The positive cases are tied to screenshots,
expected-answer checks, format constraints, and excerpt-level evidence, while
safety framing remains a useful but secondary descriptor.

\paragraph{Evidence discipline and correction.}

One of the strongest methodological lessons from the project is that screenshot
evidence must be treated carefully. Without prompt-echo cleaning and targeted
screenshot correction, the thesis would have reported two misleading claims: a
false format failure for S038 and an inflated social over-refusal story. The
corrected result is smaller but more trustworthy.

This is a general lesson for app audits. Consumer-app evidence is closer to the
real user surface, but it is harder to collect cleanly than API evidence. Screenshots
must be named, archived, transcribed, checked, and sometimes corrected. The audit
pipeline should preserve both the raw evidence and the correction audit. That is
why the appendix includes the correction table and reviewed coding summary.

\section{Cross-Vendor Context}

The Claude and Gemini API baselines show that raw API conditions differ by
vendor. Even when restricted to APIs, vendors differ on refusal-expected prompts.
This is useful context, but it is not another app-comparison result. It
reinforces the general methodological point that a single access point is not
representative of the LLM ecosystem and may not be representative of a provider's
own consumer experience.

The thesis keeps the vendor baselines secondary because there are no Claude or
Gemini app collections. Treating those API baselines as equivalent to app
comparisons would repeat the same methodological mistake the thesis criticizes.

The four-channel agreement analysis adds only secondary context. On refusal
posture, the app agrees more often with the Gemini and Claude APIs than with the
OpenAI API, because the plain OpenAI Responses API is the lone channel that
issues bare full refusals in this collection. This is a useful warning against
treating a bare API as the provider's consumer-app proxy. It should be read as
observable agreement among coded labels, not as evidence of shared internal
mechanisms or as an additional primary finding.

\section{Boundaries and Recommendations}

The thesis does not show that one deployment surface is globally better. It does
not show that the app always uses retrieval, that the API never uses safety
context, or that refusal differences are caused by a particular hidden policy.
It also does not show that the observed rates generalize to all possible prompts.
The prompt bank is diagnostic, not representative.

The correct interpretation is narrower and stronger: under documented collection
conditions, the bare API condition and productized app deployment produced
different observable behavior on a frozen benchmark-grounded prompt bank. Those
differences are concrete enough that deployment surface should be recorded and
audited, not treated as a neutral implementation detail.

\paragraph{Practical recommendations.}

For researchers, the recommendation is to report the deployment surface as part
of the experimental condition and avoid generalizing from API-only evidence to
consumer apps without validation. For benchmark designers, the recommendation is
to add fields for UI-surface signatures and tool visibility when app collection
is part of the study. For safety evaluators, the recommendation is to distinguish
bare refusal from safe redirection. For thesis-scale audits, the recommendation
is to keep a correction audit so that OCR and UI artifacts do not become
behavioral claims.

These recommendations are modest but actionable. They do not require internal
provider access. They require only disciplined black-box collection, paired
comparisons, visible-surface coding, and conservative interpretation.
